% =====================================================================
%  Leçon 201 — Espaces de fonctions. Exemples et applications.
%  Fichier autonome (préambule identique à template-lecon.tex).
%  Compilation : pdflatex (2 passes pour le nombre de pages).
% =====================================================================
\documentclass[10pt,a4paper]{article}

% ---------- Encodage, langue, fontes ----------
\usepackage[utf8]{inputenc}
\usepackage[T1]{fontenc}
\usepackage{lmodern}
\usepackage[french]{babel}

% ---------- Mise en page ----------
\usepackage[a4paper,landscape,margin=1.2cm,top=1.6cm,bottom=1.2cm,headsep=0.3cm]{geometry}
\usepackage{multicol}
\setlength{\columnsep}{1cm}
\setlength{\columnseprule}{0.4pt}
\raggedcolumns   % pas d'étirement vertical des colonnes pleines
\setlength{\parindent}{0pt}
\setlength{\parskip}{2pt}
\usepackage{lastpage}
\usepackage{fancyhdr}
\usepackage[dvipsnames]{xcolor}
\usepackage{titlesec}
\usepackage{enumitem}
\setlist{nosep,leftmargin=*}

% ---------- Mathématiques ----------
\usepackage{amsmath,amssymb,amsthm}
\usepackage{mathtools}
\usepackage{mathrsfs}
\usepackage{dsfont}

% ---------- Informations sur la leçon (à remplir) ----------
\newcommand{\numlecon}{201}
\newcommand{\titrelecon}{Espaces de fonctions. Exemples et applications.}
\newcommand{\auteur}{Grenier}   % à remplir (laisser vide pour masquer l'encart)
\newcommand{\session}{2026}   % à remplir
\newcommand{\classement}{2}   % à remplir

% ---------- En-tête / pied de page ----------
\pagestyle{fancy}
\fancyhf{}
\fancyhead[L]{\small\textbf{Leçon \numlecon}\ --- \titrelecon}
\fancyhead[R]{\small\thepage/\pageref{LastPage}}
\renewcommand{\headrulewidth}{0.4pt}

% ---------- Titres de sections (I), II)... et 1), 2)...) ----------
\renewcommand{\thesection}{\Roman{section}}
\renewcommand{\thesubsection}{\arabic{subsection}}
\titleformat{\section}{\large\bfseries}{\thesection)}{0.4em}{}
\titlespacing*{\section}{0pt}{6pt}{3pt}
\titleformat{\subsection}{\normalsize\bfseries}{\thesubsection)}{0.4em}{}
\titlespacing*{\subsection}{0pt}{4pt}{2pt}

% ---------- Environnements d'énoncés (compteur unique) ----------
\newtheoremstyle{plan}{2pt}{2pt}{}{}{\bfseries}{ :}{ }{\thmname{#1}\thmnumber{ #2}\thmnote{ (#3)}}
\newtheoremstyle{planex}{2pt}{2pt}{}{}{\bfseries\itshape}{ :}{ }{\thmname{#1}\thmnote{ (#3)}}
\theoremstyle{plan}
\newtheorem{theoreme}{Théorème}
\newtheorem{definition}[theoreme]{Définition}
\newtheorem{proposition}[theoreme]{Proposition}
\newtheorem{lemme}[theoreme]{Lemme}
\newtheorem{corollaire}[theoreme]{Corollaire}
\newtheorem{application}[theoreme]{Application}
\newtheorem{remarque}[theoreme]{Remarque}
\newtheorem{notation}[theoreme]{Notation}
\newtheorem{exemple}[theoreme]{Exemple}
\newtheorem{contreexemple}[theoreme]{Contre-exemple}
\theoremstyle{planex}   % variantes non numérotées : Ex / Rq / Contre-ex
\newtheorem*{exemple*}{Ex}
\newtheorem*{remarque*}{Rq}
\newtheorem*{contreexemple*}{Contre-ex}

% ---------- Références (équivalent de la marge du manuscrit) ----------
% \src{GOU}{31} à la fin d'un énoncé : la référence est poussée à droite
% de la dernière ligne (ou sur la ligne suivante s'il n'y a pas la place).
\newcommand{\src}[2]{\unskip\nobreak\hfill\penalty50\hskip1em\null\nobreak\hfill\mbox{\normalfont\scriptsize\textcolor{gray}{[#1~#2]}}\par}
% \srcs{ELA, GOU} dans un titre de section
\newcommand{\srcs}[1]{\ {\normalfont\small\textcolor{gray}{[#1]}}}

% ---------- Développements ----------
% \dev{1} dans le titre d'un énoncé : \begin{theoreme}[Plancherel\dev{1}]
\newcommand{\dev}[1]{\ \textcolor{red}{\textbf{[Dév.\,#1]}}}
% \begin{devbloc}{1} ... \end{devbloc} : encadre en rouge plusieurs énoncés
\newsavebox{\devsavebox}
\newenvironment{devbloc}[1]{\par\smallskip\noindent\begin{lrbox}{\devsavebox}\begin{minipage}{\dimexpr\columnwidth-2\fboxsep-2\fboxrule}\vspace{-2pt}\hfill\textcolor{red}{\textbf{Dév.\,#1}}\par\vspace{-6pt}}{\end{minipage}\end{lrbox}{\color{red}\fbox{\color{black}\usebox{\devsavebox}}}\par\smallskip}
% \listedev{...} : récapitulatif en fin de plan
\newcommand{\listedev}[1]{\par\vspace{4pt}\hrule\vspace{3pt}{\small\textcolor{red}{\textbf{Développements :}} #1}}

% ---------- Encart auteur ----------
\newcommand{\encartauteur}{\ifx\auteur\empty\else\begin{center}\fbox{\small\textbf{Auteur} : \auteur\quad$\bullet$\quad\textbf{Session} \session\quad$\bullet$\quad\textbf{Classement} : \classement}\end{center}\fi}

% ---------- Blocs de fin de plan ----------
\newcommand{\blocfin}[1]{\par\vspace{4pt}{\small\textbf{#1 :}}\par\vspace{1pt}}

% ---------- Raccourcis mathématiques ----------
\newcommand{\R}{\mathbb{R}}
\newcommand{\C}{\mathbb{C}}
\newcommand{\N}{\mathbb{N}}
\newcommand{\Z}{\mathbb{Z}}
\newcommand{\Q}{\mathbb{Q}}
\newcommand{\K}{\mathbb{K}}
\newcommand{\E}{\mathbb{E}}
\newcommand{\PP}{\mathbb{P}}
\newcommand{\T}{\mathbb{T}}
\newcommand{\U}{\mathbb{U}}
\newcommand{\Cc}{\mathcal{C}}
\newcommand{\Bb}{\mathcal{B}}
\newcommand{\Oo}{\mathcal{O}}
\newcommand{\Ll}{\mathcal{L}}
\newcommand{\Ff}{\mathcal{F}}
\newcommand{\tribu}{\mathcal{A}}
\newcommand{\ind}{\mathds{1}}
\newcommand{\norm}[1]{\left\lVert #1\right\rVert}
\newcommand{\abs}[1]{\left\lvert #1\right\rvert}
\newcommand{\ps}[2]{\left\langle #1,#2\right\rangle}
\newcommand{\Lp}[1][p]{L^{#1}}
\newcommand{\tq}{\ \text{tq}\ }
\newcommand{\ssi}{\text{ ssi }}
\DeclareMathOperator{\Supp}{Supp}
\DeclareMathOperator{\Vect}{Vect}

% =====================================================================
\begin{document}

\begin{center}
{\LARGE\bfseries Leçon \numlecon}\\[2pt]
{\large\bfseries \titrelecon}
\end{center}
\encartauteur
\vspace{2pt}\hrule\vspace{4pt}

\begin{multicols}{2}

\section{Fonctions continues sur un compact}
$(K,d)$ espace métrique compact.

\subsection{Définitions et premières propriétés}

\begin{definition}
$f : K \to (Y,d')$ métrique est continue en $x \in K$ si $\forall \varepsilon > 0,\ \exists \eta > 0 \mid \forall y \in K,\ d(x,y) \le \eta \Rightarrow d'(f(x),f(y)) \le \varepsilon$. $f$ est continue sur $K$ si elle est continue en tout $x \in K$. On note $\Cc(K,Y) = \{ f : K \to Y \mid f \text{ continue sur } K \}$.
\end{definition}

\begin{exemple*}
$K = [a,b]$, $Y = \R$ : continuité usuelle sur $\R$.
\end{exemple*}

\begin{proposition}
Si $f, g \in \Cc(K,Y)$, $E = \{ d'(f(x),g(x)),\ x \in K \}$ est borné et $\sup E$ est atteint en $x_0 \in K$. $\Cc(K,Y)$ est alors métrique pour $d_{\Cc(K,Y)}(f,g) = \sup_{x \in K} d'(f(x),g(x))$.
\end{proposition}

\begin{proposition}
Si $Y$ est complet, $(\Cc(K,Y), d_{\Cc(K,Y)})$ aussi. Si $Y$ est un Banach, $\Cc(K,Y)$ aussi pour $\norm{f}_\infty = \sup_{x \in K} \norm{f(x)}_Y$. Si $Y$ est une algèbre, $\Cc(K,Y)$ aussi.

En particulier, une limite uniforme de fonctions continues sur $K$ est continue.
\end{proposition}

\begin{exemple*}
Sur $\R$, $x \mapsto \sum_{n \ge 0} e^{-n x^2}$ est continue sur $]0,+\infty[$.
\end{exemple*}

\begin{application}[Cauchy-Lipschitz]
$E$ un espace de Banach et $F : E \to E$ tq $\norm{F(u) - F(v)} \le L \norm{u - v}$ $\forall u, v \in E$ ($L > 0$). Alors pour tout $u_0 \in E$, il existe $u \in \Cc^1([0,+\infty[, E)$ unique telle que :
\[ \begin{cases} \dfrac{du}{dt} = F(u) & \text{sur } [0,+\infty[ \\ u(0) = u_0. \end{cases} \]
\end{application}

\subsection{Théorèmes fondamentaux}

\begin{theoreme}[Heine]
$f \in \Cc(K,Y)$. $f$ est uniformément continue : $\forall \varepsilon > 0,\ \exists \eta > 0 \mid \forall x, y \in K,\ d(x,y) \le \eta \Rightarrow d'(f(x),f(y)) \le \varepsilon$.
\end{theoreme}

\begin{exemple*}
$f : [0,1[ \to \R$ uniformément continue. Alors $f$ se prolonge par continuité sur $[0,1]$.
\end{exemple*}

\begin{theoreme}[Dini]
$Y = \R$. $(f_n) \in \Cc(K,Y)^\N$ croissante. Si $(f_n)$ converge vers $f \in \Cc(K,Y)$ sur $K$ alors $f_n \to f$ dans $\Cc(K,Y)$.
\end{theoreme}

\begin{theoreme}[Arzelà]
$Y$ complet. $A \subset \Cc(K,Y)$. $\overline{A}$ est compacte ssi :
\begin{itemize}
\item $\forall x \in K$, $\{ f(x),\ f \in A \} \subset Y$ est compact.
\item Équicontinuité : $\forall \varepsilon > 0,\ \exists \eta > 0 \mid \forall f \in A,\ \forall x, y \in K,\ d(x,y) \le \eta \Rightarrow d'(f(x),f(y)) \le \varepsilon$.
\end{itemize}
\end{theoreme}

\begin{exemple*}
$M > 0$ : $\{ f : [0,1] \to \R \mid f\ M\text{-lipschitz},\ f(0) = 0 \}$ est compacte dans $\Cc([0,1],\R)$. % [correction] ajout de f(0)=0 (sinon l'ensemble n'est pas borné)
\end{exemple*}

\begin{application}
$T : \Cc([0,1],\R) \to \Cc([0,1],\R)$, $f \mapsto \left( x \mapsto \int_0^x f(s)\,ds \right)$ est un opérateur compact.
\end{application}

\begin{theoreme}[Stone-Weierstrass] % [correction] hypothèse « x ≠ y » (le manuscrit hésite avec « a ≠ b »)
$A \subset \Cc(K,\R)$ une sous-algèbre tq pour tous $x, y \in K$, $a, b \in \R$ avec $x \ne y$, il existe $f \in A$ tq $f(x) = a$, $f(y) = b$. Alors $\overline{A} = \Cc(K,\R)$.

Si $A \subset \Cc(K,\C)$ vérifie les mêmes hypothèses et si $\forall f \in A$, $\overline{f} \in A$, alors $\overline{A} = \Cc(K,\C)$.
\end{theoreme}

\begin{corollaire}[Weierstrass]
$f : [a,b] \to \R$ continue. $f$ est limite uniforme de polynômes à coefficients réels.
\end{corollaire}

\begin{exemple*}
L'ensemble des polynômes trigonométriques $\C[e^{ix}, e^{-ix}]$ est dense dans $\Cc(\U,\C)$ (Fejér).
\end{exemple*}

\begin{exemple*}
$f : [0,1] \to \R$ continue. $B_n(f) \to f$ uniformément, où
\[ B_n(f)(x) = \sum_{k=0}^{n} f\!\left(\tfrac{k}{n}\right) \binom{n}{k} x^k (1-x)^{n-k}. \]
\end{exemple*}

\section{Espaces $\Lp$}

\subsection{Généralités}

\begin{definition}
$(X, \tribu, \mu)$ mesuré. Pour $p \in [1,+\infty[$, on définit
\[ \Lp(X) = \left\{ f : X \to \C \,\Big|\, \int \abs{f}^p\,d\mu < \infty \right\} \Big/ \sim \]
où $f \sim g \Leftrightarrow f = g$ $\mu$-pp,
et $\Lp[\infty](X) = \{ f : X \to \C \mid \exists M > 0 \mid \abs{f} \le M\ \mu\text{-pp} \}$. On pose :
\begin{itemize}
\item $\norm{f}_\infty = \inf \{ M > 0 \mid \abs{f} \le M\ \mu\text{-pp} \}$ sur $\Lp[\infty](X)$,
\item $\norm{f}_p = \left( \int \abs{f}^p\,d\mu \right)^{1/p}$ sur $\Lp(X)$,
\end{itemize}
bien définis.
\end{definition}

\begin{theoreme}
$\norm{\cdot}_p$ définit une norme sur $\Lp(X)$ qui en fait un espace de Banach (Riesz-Fischer).
\end{theoreme}

\begin{corollaire}
$L^2(X)$ est un espace de Hilbert.
\end{corollaire}

\begin{application}[Radon-Nikodym]
Soit $\lambda, \mu$ deux mesures $\sigma$-finies sur $(X,\tribu)$ tq $\lambda \ll \mu$ (absolue continuité). Alors il existe une unique $h$ $\mu$-mesurable positive tq $\forall E \in \tribu$, $\lambda(E) = \int_E h\,d\mu$.
\end{application}

\begin{remarque*}
C'est encore vrai si $\lambda$ est une mesure complexe.
\end{remarque*}

\begin{application}[Espérance conditionnelle]
$(\Omega, \tribu, \PP)$ espace probabilisé. $\mathcal{G}$ une sous-tribu de $\tribu$. $X \ge 0$ ou intégrable : il existe une unique v.a. $\E(X \mid \mathcal{G})$ $\mathcal{G}$-mesurable tq $\forall A \in \mathcal{G}$, $\E(X \mathbf{1}_A) = \E\big( \E(X \mid \mathcal{G}) \mathbf{1}_A \big)$.
\end{application}

\subsection{Cas $\Omega$ ouvert de $\R^N$}
$\Omega$ ouvert de $\R^N$, les espaces se réfèrent aux fonctions sur $\Omega$. $\lambda$ est la mesure de Lebesgue. $1 \le p < \infty$.

\begin{theoreme}[densité]
$\Cc_c(\Omega) = \{ f : \Omega \to \C \mid f \text{ continue à support compact} \}$ est dense dans $\Lp(\Omega)$.
\end{theoreme}

\begin{application}
$f \in \Lp$, $y \in \R^N$. $\tau_y f = f(\cdot - y) \in \Lp$. Alors $\tau_y f \xrightarrow[y \to 0]{} f$ dans $\Lp$.
\end{application}

\begin{definition}
On appelle suite régularisante toute suite $(\rho_n)_{n \ge 0}$ tq $\rho_n \in \Cc_c^\infty(\R^N)$, $\Supp \rho_n \subset B(0, \frac{1}{n})$, $\int \rho_n = 1$, $\rho_n \ge 0$.
\end{definition}

\begin{remarque*}
On pose $\rho(x) = \begin{cases} e^{-\frac{1}{1 - \norm{x}^2}} & \text{si } \norm{x} < 1 \\ 0 & \text{si } \norm{x} \ge 1 \end{cases}$ et $\rho_n(x) = C n^N \rho(nx)$, $C = \left( \int \rho \right)^{-1}$ est une suite régularisante.
\end{remarque*}

\begin{theoreme}
Si $f \in \Lp(\R^N)$, $\rho_n * f \xrightarrow[n \to \infty]{} f$ dans $\Lp(\R^N)$ où $\rho_n * f : x \mapsto \int f(y) \rho_n(x-y)\,dy$.
\end{theoreme}

\begin{corollaire}
$\Omega \subset \R^N$ ouvert. $\Cc_c^\infty(\Omega)$ dense dans $\Lp(\Omega)$.
\end{corollaire}

\begin{exemple*}
$f \in L^1(\R)$, $\hat{f}(n) = \int f(x) e^{inx}\,dx \xrightarrow[n \to \infty]{} 0$ (lemme de Riemann-Lebesgue).
\end{exemple*}

\subsection{Applications en analyse fonctionnelle}

\begin{theoreme}[Riesz]
$1 < p < \infty$. $\Omega$ ouvert de $\R^N$. $\varphi$ forme linéaire sur $\Lp_\R(\Omega)$, continue : $\varphi \in (\Lp)'$. Alors $\exists!\ u \in L^q$ tq $\forall f \in \Lp$, $\varphi(f) = \int u f\,d\lambda$ ($\frac{1}{p} + \frac{1}{q} = 1$). $(\Lp)'$ est isométrique à $L^q$.
\end{theoreme}

\begin{theoreme}[Plancherel\dev{1}]
La transformation de Fourier réalise une isométrie de $L^2(\R) \cap L^1(\R)$ dans $L^2$ pour $\norm{\cdot}_2$, donc se prolonge en une isométrie de $L^2(\R)$. Cette isométrie est bijective.
\end{theoreme}

\begin{definition}
\[ L^2(\T) = \Big\{ f : \R \to \C,\ 2\pi\text{-pér. } \lambda\text{-pp tq } \textstyle\int_0^{2\pi} \abs{f}^2\,d\lambda < \infty \Big\} \]
et $\norm{f}_2 = \left( \frac{1}{2\pi} \int_0^{2\pi} \abs{f}^2\,d\lambda \right)^{1/2}$.
C'est un Hilbert.
\end{definition}

\begin{theoreme}
Pour $n \in \Z$, $e_n : x \mapsto e^{inx} \in L^2(\T)$. $(e_n)_{n \in \Z}$ forme une base hilbertienne de $L^2(\T)$.
\end{theoreme}

\begin{corollaire}
Si $u \in L^2(\T)$, la série de Fourier $(S_u)$ définie par $S_u = \sum_{n \in \Z} \ps{e_n}{u} e_n$ converge vers $u$ dans $L^2(\T)$. $u = \sum_{n \in \Z} \ps{e_n}{u} e_n$ et $\norm{u}_2^2 = \sum_{n \in \Z} \abs{\ps{e_n}{u}}^2$.
\end{corollaire}

\begin{exemple*}
$\displaystyle \int_\R \left( \frac{\sin x}{x} \right)^2 dx = \pi$, $\displaystyle \sum_{n=1}^{\infty} \frac{1}{n^4} = \frac{\pi^4}{90}$.
\end{exemple*}

\section{Espaces de fonctions holomorphes}
$\Omega \subset \C$ ouvert. On note $\Oo(\Omega)$ l'ensemble des fonctions holomorphes sur $\Omega$. Soit $(K_n)_{n \ge 0}$ une suite exhaustive de compacts de $\Omega$ ($\bigcup K_n = \Omega$, $K_n \subset \mathring{K}_{n+1}$).

\begin{definition} % [correction] f, g ∈ O(Ω) (le manuscrit écrit f, g ∈ Ω)
Si $f, g \in \Oo(\Omega)$, on pose $p_n(f) = \sup_{y \in K_n} \abs{f(y)}$ et $d(f,g) = \sum_{i \ge 1} 2^{-i} \min(1, p_i(f-g))$. C'est une distance sur $\Oo(\Omega)$ et $f_n \xrightarrow{d} f$ ssi $f_n \to f$ uniformément sur tout compact de $\Omega$.
\end{definition}

\begin{theoreme}[Weierstrass]
Si $(f_n) \in \Oo(\Omega)^\N$, $f : \Omega \to \C$ et $f_n \to f$ uniformément sur tout compact de $\Omega$, alors $f \in \Oo(\Omega)$ et $\forall k \ge 0$, $f_n^{(k)} \to f^{(k)}$ dans $\Oo(\Omega)$.
\end{theoreme}

\begin{remarque}
$(\Oo(\Omega), d)$ est complet. C'est un espace stable par dérivation. Il n'est pas normable.
\end{remarque}

\begin{theoreme}[Montel\dev{2}]
$H \subset \Oo(\Omega)$. Alors $H$ est relativement compact ssi pour tout compact $K \subset \Omega$, la famille $(f|_K)_{f \in H}$ est bornée pour la distance uniforme. Une telle famille est dite normale.
\end{theoreme}

\end{multicols}

\listedev{1. Théorème de Plancherel (Thm~22). --- 2. Théorème de Montel (Thm~29).}

\end{document}
